% Generated by Listas/make_statements.py from pratica-lecture-06.md, item 9. Do not edit.
\textbf{Source.} Tadelis (2013), \emph{Game Theory: An Introduction}, Exercise 10.8 ``Tacit Collusion'', p.~218 {[}PDF 237{]}. Reproduced verbatim. Answer in the source: none in the book (the preface, p.~xiii, says about half are solved on Princeton UP's website; not checked).



\begin{quote}

\textbf{Tacit Collusion:} Two firms, which have zero marginal cost and no fixed cost, produce some good, each producing \(q_i \ge 0\), \(i \in \{1, 2\}\). The demand for this good is given by \(p = 200 - Q\), where \(Q = q_1 + q_2\).



\begin{enumerate}

\def\labelenumi{\alph{enumi}.}

\item

  First consider the case of Cournot competition, in which each form chooses \(q_i\) and this game is infinitely repeated with a discount factor \(\delta < 1\). Solve for the static stage-game Cournot-Nash equilibrium.

\item

  For which values of \(\delta\) can you support the firms' equally splitting monopoly profits in each period as a subgame-perfect equilibrium that uses grim-trigger strategies (i.e., after one deviates from the proposed split, they resort to the static Cournot-Nash equilibrium thereafter)? (Note: Be careful in defining the strategies of the firms!)

\item

  Now assume that the firms compete à la Bertrand, each choosing a price \(p_i \ge 0\), where the lowest-price firm gets all the demand and in case of a tie they split the market. Solve for the static stage-game Bertrand-Nash equilibrium.

\item

  For which values of \(\delta\) can you support the firms' equally splitting monopoly profits in each period as a subgame-perfect equilibrium that uses grim-trigger strategies (i.e., after one deviates from the proposed split, they resort to the static Bertrand-Nash equilibrium thereafter)? (Note: Be careful in defining the strategies of the firms!)

\item

  Now instead of using grim-trigger strategies, try to support the firms' equally splitting monopoly profits as a subgame-perfect equilibrium in which after a deviation firms would resort to the static Bertrand competition for only two periods. For which values of \(\delta\) will this work? Why is this answer different than your answer in (d)?

\end{enumerate}

\end{quote}

